\SourceLeaf{406}%
\par\medskip\Needspace{6\baselineskip}
\pdfbookmark[2]{VII.}{carnot-zusatz-vii}
\centerline{\textbf{VII.}}\nobreak\smallskip
\begin{center}\textit{Herleitung der Grundformeln der sphärischen Trigonometrie.}\end{center}
\hspace*{\parindent}Sind \(a,b,c\) die Seiten und \(A,B,C\) die gegenüberliegenden Winkel eines sphärischen Dreiecks, so gilt
\[
\cos a=\cos b.\cos c+\sin b.\sin c.\cos A
\]
wie \textsc{Lagrange} für den Fall, dass sowohl \(b\) als auch \(c\) kleiner als \(90^0\) sind, elegant bewiesen hat. Der Beweis lässt sich jedoch leicht auf alle anderen Fälle ausdehnen.

Die folgenden Fälle sind möglich.
\[
\text{I.}\quad .\ .\ .\ .\ .\ .\qquad b<90^0,\quad c<90^0
\]
\par
Hier gilt der Beweis unmittelbar.
\[
\text{II.}\quad .\ .\ .\ .\ .\ .\qquad b>90^0,\quad c>90^0
\]
\par
Man verlängere die Seiten \(b\) und \(c\) über die Punkte \(B\) und \(C\) hinaus bis zu ihrem Schnittpunkt \(A'\) und bestimme den Wert von \(a\) anhand des Dreiecks \(A'BC\).
\[
\text{III.}\quad .\ .\ .\ .\ .\ .\qquad b>90^0,\quad c<90^0
\]
\par
Man verlängere die Seiten \(b\) und \(a\) über die Punkte \(A\) und \(B\) hinaus bis zu ihrem Schnittpunkt \(C'\). Im Dreieck \(C'AB\) folgt dann aus Fall I
\[
\cos(180^0-a)=\cos c.\cos(180^0-b)+\sin c.\sin(180^0-b).\cos(180^0-A)
\]
Das ist mit der Grundformel identisch.

Der Fall \(b<90^0\) und \(c>90^0\) ist im Wesentlichen derselbe.
\[
\text{IV.}\quad .\ .\ .\ .\ .\ .\qquad b=90^0,\quad A=90^0
\]
\par
Hier ist \(C\) der Pol von \(AB\), also gilt notwendig auch \(a=90^0\). Damit ist die Formel
\[
\cos a=\cos b.\cos c+\sin b.\sin c.\cos A
\]
unmittelbar ersichtlich. Der Fall \(c=90^0,\ A=90^0\) ist im Wesentlichen derselbe.
